\chapter{LesFurets.com}

\section{L'assurance}
Le secteur de l’assurance fait partie du secteur économique tertiaire qui regroupe les industries de services essentiellement immatériels. Il s’agit d’un secteur d’activité où les acteurs économiques vendent une protection contre le risque avec un poids très important au sein du paysage économique mondial et, plus particulièrement, au sein du paysage économique français. Sa particularité réside dans son activité et sa diversité, puisqu'il rassemble une multitude d’acteurs : compagnies d’assurance, courtiers, mutuelles, acteurs internet, etc... Mondialement, le marché de l’assurance est actuellement dominé par les états-Unis suivi par le Japon. La France occupe la 4ème position, derrière le Royaume Uni, avec 6,5\% des cotisations mondiales, soit près de 211 235 millions d’euros. (FFSA et Gema)

\section{Les comparateurs d'assurance} Un comparateur de prime d’assurance est un outil informatique permettant de comparer toutes les offres d’assurances répondant à un ensemble de critères mentionnés par le client. Avec un marché fortement concurrentiel et très complexe, les comparateurs ont un rôle de plus en plus important auprès des clients en matière de comparaison, de choix et de souscription de produits d’assurance. Aujourd’hui, le marché est occupé par des concurrents connus tels que LeLynx.fr, LesFurets.com, Assurland.com, Hyperassur, automotocompare.fr, kelassur.com, etc... En effet, ces comparateurs présentent des environnements ergonomiques permettant aux visiteurs d’effectuer un comparatif de contrats auto, moto, multirisques habitation, emprunteur ou santé, et de sélectionner les offres les mieux adaptées à leurs attentes. Pour répondre aux besoins des clients, chaque acteur maintient son propre modèle qui repose sur plusieurs aspects : le tarif, la qualité des services proposés, les services après-vente, les garanties, les franchises, les prestations, etc. Ainsi, chaque comparateur présente des offres d’assureurs différents. Le référencement web et la publicité présentent des investissements lourds pour l’ensemble des concurrents qui cherchent à obtenir une visibilité conséquente auprès de leur clients. Les nouveaux arrivants doivent forcément investir massivement pour trouver une place parmi les acteurs historiques sur ce marché.

\section{Historique de l'entreprise}
Courtanet, fondée en 2005 par Jehan de Castet, est la société éditrice du site AssureMieux.com. Avec ses sites www.assuremieux.com, www.creditmieux.com, et son logiciel Bénéfit, la société est devenue rapidement l’un des plus importants fournisseurs français de solutions de comparaison d’assurances. Plus de 1 500 courtiers indépendants français étaient équipés de ses solutions. Depuis Novembre 2010, Courtanet est majoritairement détenu par BGL Group créé en 1992, intermédiaire britannique d’assurances et propriétaire également du leader de la comparaison d’assurances en Grande Bretagne via www.comparethemarket.com, qui a emporté 35\% du marché au Royaume-Uni. Le groupe est aujourd’hui structuré autour de quatre «piliers» : Entreprises intermédiaires, Entreprises pilotées par marque comme LesFurets pour la France, HoyHoy pour la Hollande, Beagle Street pour le Royaume Uni et d'autre services transverses. En Avril 2012, la plateforme de comparaison Assuremieux change d’identité et dévoile une nouvelle marque baptisée LesFurets.com, le comparateur qui « simplifie » l’assurance. LesFurets.com est un comparateur indépendant et impartial d’assurance auto, moto, santé, habitation et crédit. Sa mission est d’aider les consommateurs à trouver l’assurance la mieux adaptée à leurs besoins au meilleur tarif. La mission du site est d’apporter plus de transparence et de simplicité dans l’assurance. Ainsi, il permet de comparer facilement les tarifs, les garanties, les franchises et les services de grands assureurs auto, santé,moto, habitation et crédit. Son objectif est d’aider à trouver l’assurance la mieux adaptée aux besoins des clients au meilleur tarif. Dans le but d’offrir un service de qualité, LesFurets.com s’associe avec les plus grands assureurs du marché : Direct Assurance, Amaguiz, Allainz, Eurofil, AllSecur, Aon Asurances, AcommeAssure, SOS Malus, Assurpeople, L’olivier Assurances, 4assur, EuroAssurance, ActiveAssurances, AssurOnline, assurbike, aloa Assurances, SwissLife, Alptis, etc.
\begin{center}
\vspace{2cm}
\newlength{\hauteur}
\setlength{\hauteur}{3cm}
\includegraphics{bgl/bgl.jpg}\\
\includegraphics[height=\hauteur]{bgl/comparethemarket.jpg}
\includegraphics[height=\hauteur]{bgl/hoyhoy.png}
\includegraphics[height=\hauteur]{bgl/lesfurets.png}
\end{center}
\clearpage
\section{L'application Web}
Le site LesFurets.com est une application Web écrite en Java SE et mise à jour régulièrement avec la dernière version de Java, aujourd'hui elle est pleinement compatible avec Java 8. Elle est déployée à l'aide de Apache Maven. Apache Maven est un outil pour la gestion et l'automatisation de production des projets logiciels Java. L'objectif recherché est comparable au système Make sous Unix : produire un logiciel à partir de ses sources, en optimisant les tâches réalisées à cette fin et en garantissant le bon ordre de fabrication.
Il est semblable à l'outil Ant, mais fournit des moyens de configuration plus simples, eux aussi basés sur le format XML. Maven est géré par l'organisation Apache Software Foundation. Précédemment Maven était une branche de l'organisation Jakarta Project.
Maven utilise un paradigme connu sous le nom de Project Object Model (POM) afin de décrire un projet logiciel, ses dépendances avec des modules externes et l'ordre à suivre pour sa production. Il est livré avec un grand nombre de tâches pré-définies, comme la compilation de code Java ou encore sa modularisation.
Un élément clé et relativement spécifique de Maven est son aptitude à fonctionner en réseau. Une des motivations historiques de cet outil est de fournir un moyen de synchroniser des projets indépendants : publication standardisée d'information, distribution automatique de modules JAR. Ainsi en version de base, Maven peut dynamiquement télécharger des dépendances sur des dépôts logiciels connus. Il propose ainsi la synchronisation transparente de modules nécessaires. C'est ainsi qu'a chaque packaging, la gestion des librairies nécessaires, la compilation et l'exécution des tests présent dans le projet sont automatiquement effectuées. Enfin la base de code est mise à jour à l'aide du gestionnaire de version GIT héberger sur un serveur de l'entreprise.